% Source: Wald GR, physical PDF pp. 472--482
%         (printed pp. 459--469).
% Coverage: Appendix E.2, Hamiltonian Formulation; equations (E.2.1)--(E.2.47).
% Figures/tables/footnotes: no figures/tables; footnotes 2--3.
% Status: source-reviewed and compile-clean.
% Uncertainties: none.

\subsection{Hamiltonian Formulation}\label{sec:E.2}

A Lagrangian formulation of a field theory is "spacetime covariant." One specifies
on the spacetime manifold an action functional of the field \(\psi\) whose
extremization yields the field equations. On the other hand, a Hamiltonian
formulation of a field theory requires a breakup of spacetime into space and time.
Indeed, the first step in producing a Hamiltonian formulation of a field theory
consists of choosing a time function \(t\) and a vector field \(t^a\) on a spacetime
such that the surfaces, \(\Sigma_t\), of constant \(t\) are spacelike Cauchy surfaces
and such that \(t^a\nabla_a t=1\). The vector field \(t^a\) may be interpreted as
describing the "flow of time" in the spacetime and can be used to identify each
\(\Sigma_t\) with the initial surface \(\Sigma_0\). In Minkowski spacetime the choice
of \(t\) and \(t^a\) usually is made via a global inertial coordinate system, but in
curved spacetime there may not be any preferred choice. In performing integrals of
functions over \(M\) it would be natural for most purposes to use the volume element
\(\epsilon_{abcd}\) associated with the spacetime metric. Similarly, in performing
integrals over \(\Sigma_t\), it would be natural in most cases to use the volume
element \({}^{(3)}\epsilon_{abc}=e_{dabc}n^d\), where \(n^d\) is the unit normal to
\(\Sigma_t\). However, these volume elements will, in general, be "time dependent"
in the sense that \(\mathcal L_t e_{abcd}\ne0\) and
\(\mathcal L_t{}^{(3)}\epsilon_{abc}\ne0\). The use of a time-dependent volume
element on \(\Sigma_t\) is particularly inconvenient if we wish to identify
\(\Sigma_t\) with \(\Sigma_0\) in order to view dynamical evolution as the change of
fields on the fixed manifold \(\Sigma_0\). Therefore, we shall introduce a fixed
volume element \(e_{abcd}\) on \(M\) satisfying \(\mathcal L_t e_{abcd}=0\). [One
way to do this---at least locally---would be to introduce coordinates
\(x^1,x^2,x^3\) in addition to \(t\) such that
\(t^a=(\partial/\partial t)^a\) and to take \(e\) to be the coordinate volume
element \(dt\wedge dx^1\wedge dx^2\wedge dx^3\).] On each \(\Sigma_t\), we define
\({}^{(3)}e_{abc}=e_{dabc}t^d\). Unless otherwise stated, all integrals over \(M\)
will be performed using the volume element \(e_{abcd}\) and all integrals over
\(\Sigma_t\) will be with respect to the volume element \({}^{(3)}e_{abc}\). Thus,
in particular, in order that our results be independent of our choice of
\(e_{abcd}\), the Lagrangian density must be a scalar density on \(M\) and the
momentum \(\pi\) (defined below) must be a tensor density on \(\Sigma_t\). As
previously noted in our discussion of the Einstein Lagrangian, the introduction of
\(e_{abcd}\) could be avoided by incorporating the volume element into the
definition of \(\mathcal L\), \(\pi\), and other quantities, but we choose not to do
so since this procedure is rather cumbersome.

The next step in giving a Hamiltonian formulation is to define a configuration
space for the field by specifying what tensor field (or fields) \(q\) on \(\Sigma_t\)
physically describes the instantaneous configuration of the field \(\psi\). Let
\(\mathcal Q\) denote this configuration space. The space of possible momenta of the
field at a given configuration \(q\) then is taken to be the "cotangent space,"
\(T_q^\vee\mathcal Q\), of the configuration space at \(q\). Since the set of possible
configurations of the field is infinite-dimensional, we shall not attempt here to
give a precise definition of \(T_q^\vee\mathcal Q\). However, in the case where the
allowed infinitesimal variations (i.e., "tangent vectors") \(\delta q\) at \(q\)
are represented by tensor fields on \(\Sigma_t\) of type \((k,l)\), we shall take the
space of momenta to consist of tensor fields, \(\pi\), of type \((l,k)\) on
\(\Sigma_t\), so that \(\pi\) maps \(\delta q\) into \(\mathbb R\) via
\(\delta q\mapsto\int_{\Sigma_t}\pi\delta q\), where contraction of indices is
understood. A prescription must then be given for associating a momentum \(\pi\) to
the field \(\psi\) on \(\Sigma_t\). The final and most nontrivial step required for a
Hamiltonian formulation of a field theory is the specification of a functional
\(H[q,\pi]\) on \(\Sigma_t\), called the \emph{Hamiltonian}, which is of the form
\begin{equation}
  H=\int_{\Sigma_t}\mathcal H,
  \tag{E.2.1}\label{eq:E.2.1}
\end{equation}
where the \emph{Hamiltonian density} \(\mathcal H\) is the local function of
\(q\), \(\pi\), and of their spatial derivatives up to a finite order, such that the
pair of equations,
\begin{equation}
  \dot q\equiv\mathcal L_tq=\frac{\delta H}{\delta\pi},
  \tag{E.2.2}\label{eq:E.2.2}
\end{equation}
\begin{equation}
  \dot\pi\equiv\mathcal L_t\pi=-\frac{\delta H}{\delta q},
  \tag{E.2.3}\label{eq:E.2.3}
\end{equation}
is equivalent to the field equation satisfied by \(\psi\).

Given a Lagrangian formulation of a field theory, there is a standard prescription
for obtaining a Hamiltonian formulation which is closely analogous to the well
known procedure in ordinary particle mechanics. First, one takes \(q\) to be simply
the field \(\psi\) evaluated on \(\Sigma_t\). Then one views the Lagrangian density
as a function of \(q\), its time derivatives, and its space derivatives. Assuming that
\(\mathcal L\) does not depend on time derivatives of \(q\) higher than first order,
we take the momentum, \(\pi\), associated with \(\psi\) on \(\Sigma_t\), to be
\begin{equation}
  \pi=\frac{\partial\mathcal L}{\partial\dot q}.
  \tag{E.2.4}\label{eq:E.2.4}
\end{equation}
Next, we attempt to solve Eq.~\eqref{eq:E.2.4} for \(\dot q\) as a function of
\(q\) and \(\pi\). If this can be done, we define
\begin{equation}
  \mathcal H(q,\pi)=\pi\dot q-\mathcal L,
  \tag{E.2.5}\label{eq:E.2.5}
\end{equation}
where \(\dot q=\dot q(q,\pi)\) is understood in this equation both in its explicit
appearance and in its implicit appearance in \(\mathcal L\). With this choice of
\(\mathcal H\), Eqs.~\eqref{eq:E.2.2} and \eqref{eq:E.2.3} are equivalent to
Eq.~\eqref{eq:E.1.5}. To see this, we define
\begin{equation}
  J=\int_{t_1}^{t_2}H\,dt
    =\int_{t_1}^{t_2}dt\int_{\Sigma_t}\mathcal H
    =-S+\int_{t_1}^{t_2}dt\int_{\Sigma_t}\pi\dot q .
  \tag{E.2.6}\label{eq:E.2.6}
\end{equation}
Then, for a smooth one-parameter variation of \(\psi\) which satisfies
\(\delta\psi=0\) at \(t=t_1\) and \(t=t_2\), we have
\begin{align}
  \frac{dJ}{d\lambda}
    &=\int_{t_1}^{t_2}dt\int_{\Sigma_t}
      \left[\frac{\delta H}{\delta q}\delta q
        +\frac{\delta H}{\delta\pi}\delta\pi\right]\notag\\
    &=\int_{t_1}^{t_2}dt\int_{\Sigma_t}
      [\pi\delta\dot q+\dot q\delta\pi]-\frac{dS}{d\lambda}\notag\\
    &=\int_{t_1}^{t_2}dt\int_{\Sigma_t}
      [-\dot\pi\delta q+\dot q\delta\pi]-\frac{dS}{d\lambda},
  \tag{E.2.7}\label{eq:E.2.7}
\end{align}
where an integration by parts was performed in the last line. Thus, comparing the
first and last line of this equation, we see that \(\delta S/\delta\psi=0\) if and
only if Eqs.~\eqref{eq:E.2.2} and \eqref{eq:E.2.3} are satisfied. Thus,
\(\mathcal H\) is a Hamiltonian density for \(\psi\).

The above procedure yields in a straightforward manner a Hamiltonian formulation
of the theory of a Klein--Gordon field in Minkowski spacetime. We choose a global
inertial coordinate system to obtain \(t\) and \(t^a\) and choose \(e_{abcd}\) to be
the natural volume element \(\epsilon_{abcd}\) since
\(\mathcal L_t e_{abcd}=0\). We choose \(q\) on \(\Sigma_t\) to be \(\phi\)
evaluated on \(\Sigma_t\) and write \(\mathcal L_{\mathrm{KG}}\),
Eq.~\eqref{eq:E.1.6}, as
\begin{equation}
  \mathcal L_{\mathrm{KG}}
    =\frac{1}{2}\bigl(\dot\phi^2-\boldsymbol{\nabla}\phi\cdot\boldsymbol{\nabla}\phi
      -m^2\phi^2\bigr),
  \tag{E.2.8}\label{eq:E.2.8}
\end{equation}
where we use ordinary three-dimensional vector notation on \(\Sigma_t\). We find
\begin{equation}
  \pi=\frac{\partial\mathcal L_{\mathrm{KG}}}{\partial\dot\phi}=\dot\phi.
  \tag{E.2.9}\label{eq:E.2.9}
\end{equation}
Hence, we define the Hamiltonian density by
\begin{equation}
  \mathcal H_{\mathrm{KG}}
    =\pi\dot\phi-\mathcal L_{\mathrm{KG}}
    =\frac{1}{2}\bigl(\pi^2+\boldsymbol{\nabla}\phi\cdot\boldsymbol{\nabla}\phi+m^2\phi^2\bigr).
  \tag{E.2.10}\label{eq:E.2.10}
\end{equation}
It then may be verified that for
\(H_{\mathrm{KG}}=\int_{\Sigma_t}\mathcal H_{\mathrm{KG}}\),
Eqs.~\eqref{eq:E.2.2} and \eqref{eq:E.2.3} indeed are equivalent to the
Klein--Gordon equation. Note also that the numerical value of \(H_{\mathrm{KG}}\)
is just the total energy of the Klein--Gordon field.

For the case of the electromagnetic field in Minkowski spacetime, it is not as
straightforward to obtain a Hamiltonian formulation by this procedure. We
provisionally take \(q\) to be the vector potential \(A_a\) evaluated on \(\Sigma_t\)
and decompose it into its normal and tangential parts,
\begin{equation}
  V=-A_an^a,
  \tag{E.2.11}\label{eq:E.2.11}
\end{equation}
\begin{equation}
  {}^{(3)}A_a=h_a{}^bA_b,
  \tag{E.2.12}\label{eq:E.2.12}
\end{equation}
where \(n^a\) is the unit normal to \(\Sigma_t\) and
\(h_{ab}=\eta_{ab}+n_an_b\) is the induced spatial metric on \(\Sigma_t\). In
ordinary three-dimensional vector notation, the Lagrangian density,
Eq.~\eqref{eq:E.1.9}, is
\begin{equation}
  \mathcal L_{\mathrm{EM}}
    =\frac{1}{2}(\dot{\mathbf{A}}+\boldsymbol{\nabla} V)\cdot(\dot{\mathbf{A}}+\boldsymbol{\nabla} V)
      -\frac{1}{2}(\boldsymbol{\nabla}\times\mathbf{A})\cdot(\boldsymbol{\nabla}\times\mathbf{A}).
  \tag{E.2.13}\label{eq:E.2.13}
\end{equation}
Hence the momentum conjugate to \(\mathbf{A}\) is
\begin{equation}
  \boldsymbol{\pi}=\dot{\mathbf{A}}+\boldsymbol{\nabla} V=-\mathbf{E}.
  \tag{E.2.14}\label{eq:E.2.14}
\end{equation}
However, \(\dot V\) does not appear in \(\mathcal L_{\mathrm{EM}}\), so the momentum
\(\pi_V\) conjugate to \(V\) vanishes identically,
\begin{equation}
  \pi_V=0.
  \tag{E.2.15}\label{eq:E.2.15}
\end{equation}

Thus, we do not obtain an invertible relation between \(\pi\) and \(\dot q\).
Consequently, if we define \(\mathcal H\) by Eq.~\eqref{eq:E.2.5}, we will not be
able to eliminate \(\dot q\) in favor of \(\pi\) and \(q\), and our general
prescription for obtaining a Hamiltonian formulation breaks down. This difficulty is
directly related to the fact that there is gauge arbitrariness in \(A_a\), and hence
we cannot expect to get deterministic dynamics for \(A_a\) of the form
\eqref{eq:E.2.2}, \eqref{eq:E.2.3}.

However, this difficulty can be resolved by the following considerations. The fact
that \(\pi_V\) vanishes identically suggests that we should not view \(V\) as a
dynamical variable. It suggests that we should take the configuration field \(q\) to
be simply \(\mathbf{A}\). Therefore, we define \(\mathcal H_{\mathrm{EM}}\) by
\begin{align}
  \mathcal H_{\mathrm{EM}}
    &=\boldsymbol{\pi}\cdot\dot{\mathbf{A}}-\mathcal L_{\mathrm{EM}}\notag\\
    &=\frac{1}{2}\boldsymbol{\pi}\cdot\boldsymbol{\pi}
      +\frac{1}{2}\mathbf{B}\cdot\mathbf{B}-\boldsymbol{\pi}\cdot\boldsymbol{\nabla} V\notag\\
    &=\frac{1}{2}\boldsymbol{\pi}\cdot\boldsymbol{\pi}
      +\frac{1}{2}\mathbf{B}\cdot\mathbf{B}
      +V\boldsymbol{\nabla}\cdot\boldsymbol{\pi}-\boldsymbol{\nabla}\cdot(V\boldsymbol{\pi}),
  \tag{E.2.16}\label{eq:E.2.16}
\end{align}
where \(\mathbf{B}\equiv\boldsymbol{\nabla}\times\mathbf{A}\). The last term on the right-hand side
of Eq.~\eqref{eq:E.2.16} is a total divergence and thus contributes only a boundary
term to \(H_{\mathrm{EM}}=\int_{\Sigma_t}\mathcal H_{\mathrm{EM}}\) which vanishes
in the limit as the boundary goes to infinity for the asymptotic conditions usually
imposed on \(V\) and \(\boldsymbol{\pi}=-\mathbf{E}\). Hence we shall discard this term.

We now view \(H_{\mathrm{EM}}\) as a functional of \(\mathbf{A}\) and \(\boldsymbol{\pi}\), with
\(V\) effectively playing the role of a Lagrange multiplier, i.e., we append the
equation
\begin{equation}
  \frac{\delta H_{\mathrm{EM}}}{\delta V}=0
  \tag{E.2.17}\label{eq:E.2.17}
\end{equation}
to Eqs.~\eqref{eq:E.2.2} and \eqref{eq:E.2.3} for \(\mathbf{A}\) and \(\boldsymbol{\pi}\).
Equation \eqref{eq:E.2.17} yields
\begin{equation}
  \boldsymbol{\nabla}\cdot\mathbf{E}=0,
  \tag{E.2.18}\label{eq:E.2.18}
\end{equation}
whereas Eqs.~\eqref{eq:E.2.2} and \eqref{eq:E.2.3} yield, respectively,
\begin{equation}
  \dot{\mathbf{A}}=\frac{\delta H_{\mathrm{EM}}}{\delta\boldsymbol{\pi}}
    =\boldsymbol{\pi}-\boldsymbol{\nabla} V=-\mathbf{E}-\boldsymbol{\nabla} V,
  \tag{E.2.19}\label{eq:E.2.19}
\end{equation}
\begin{equation}
  \dot{\boldsymbol{\pi}}=-\dot{\mathbf{E}}
    =-\frac{\delta H_{\mathrm{EM}}}{\delta\mathbf{A}}
    =-\boldsymbol{\nabla}\times(\boldsymbol{\nabla}\times\mathbf{A}).
  \tag{E.2.20}\label{eq:E.2.20}
\end{equation}
Thus, we see that the system of Eqs.~\eqref{eq:E.2.18}--\eqref{eq:E.2.20} is
equivalent to Maxwell's equations. Furthermore, we obtain from this formulation a
natural breakup of Maxwell's equations into the constraint \eqref{eq:E.2.18} and
the evolution equations \eqref{eq:E.2.19}, \eqref{eq:E.2.20}. Note that, again,
the numerical value of \(H_{\mathrm{EM}}\) for a solution of Maxwell's equations is
proportional to the total energy of the electromagnetic field.

Thus, we have obtained a Hamiltonian formulation of Maxwell's equations in
Minkowski spacetime which has the feature that a non-dynamical variable appears in
\(\mathcal H\) and effectively plays the role of a Lagrange multiplier enforcing the
constraint \eqref{eq:E.2.17}. This type of Hamiltonian formulation is called a
\emph{constrained Hamiltonian formulation}. As will be discussed further below, it
can be expected to arise in any theory where the field variables have a gauge
arbitrariness.

We turn, now, to the task of obtaining a Hamiltonian formulation of Einstein's
equations. As in the previous cases, we choose a time function \(t\) and a "time
flow" vector field \(t^a\) on \(M\) satisfying \(t^a\nabla_at=1\). Note, however,
that in this case one cannot interpret \(t\) and \(t^a\) in terms of physical
measurements using clocks until one knows the spacetime metric, which, of course,
is the unknown field variable in Einstein's equation. Given a metric \(g_{ab}\), it
is convenient to decompose \(t^a\) into its normal and tangential parts with respect
to the surfaces, \(\Sigma_t\), of constant \(t\). As in chapter 10, we define the
\emph{lapse function}, \(N\), by
\begin{equation}
  N=-g_{ab}t^an^b=(n^a\nabla_at)^{-1},
  \tag{E.2.21}\label{eq:E.2.21}
\end{equation}
and the \emph{shift vector} \(N^a\) by
\begin{equation}
  N^a=h^a{}_bt^b,
  \tag{E.2.22}\label{eq:E.2.22}
\end{equation}
where again \(n^a\) is the unit normal to \(\Sigma_t\) and
\(h_{ab}=g_{ab}+n_an_b\) is the induced spatial metric on \(\Sigma_t\). Thus,
\(N\) measures the rate of flow of proper time, \(\tau\), with respect to coordinate
time, \(t\), as one moves normally to \(\Sigma_t\), whereas \(N^a\) measures the
amount of "shift" tangential to \(\Sigma_t\) contained in the time flow vector field
\(t^a\) (see Fig.~10.2 of chapter 10). In terms of \(N\), \(N^a\), and \(t^a\), we
have
\begin{equation}
  n^a=\frac{1}{N}(t^a-N^a),
  \tag{E.2.23}\label{eq:E.2.23}
\end{equation}
and hence the inverse spacetime metric can be written as
\begin{equation}
  g^{ab}=h^{ab}-n^an^b
    =h^{ab}-N^{-2}(t^a-N^a)(t^b-N^b).
  \tag{E.2.24}\label{eq:E.2.24}
\end{equation}
It is convenient to choose as our field variables the spatial metric, \(h_{ab}\), the
lapse function \(N\), and the covariant form of the shift vector,
\(N_a=h_{ab}N^b\) rather than the inverse metric, \(g^{ab}\), which was used as the
field variable in the previous section. The requirements that \(h^{ac}h_{cb}\) be
the identity operator on the tangent space to \(\Sigma_t\) and that
\(h^{ab}\nabla_bt=0\) allow us to compute \(h^{ab}\) from \(h_{ab}\) and thence
obtain \(N^a=h^{ab}N_b\). Thus, from Eq.~\eqref{eq:E.2.24} we see that the
information contained in \((h_{ab},N,N_a)\) is equivalent to that contained in
\(g^{ab}\).

Again, we shall use a fixed volume element \(e_{abcd}\) on spacetime satisfying
\(\mathcal L_t e_{abcd}=0\) and will use the volume element
\({}^{(3)}e_{abc}=e_{dabc}t^d\) on \(\Sigma_t\). We note in analogy with the
remarks below Eq.~\eqref{eq:E.1.10}, we have
\({}^{(3)}\epsilon_{abc}=\sqrt{h}\,{}^{(3)}e_{abc}\), where \(h\) is the
determinant of the matrix of components, \(h_{\mu\nu}\), of \(h_{ab}\) in a basis
where the nonvanishing components of \({}^{(3)}e_{abc}\) have the values \(\pm1\).
It then follows that
\begin{equation}
  \sqrt{-g}=N\sqrt{h}.
  \tag{E.2.25}\label{eq:E.2.25}
\end{equation}
The first step in obtaining a Hamiltonian functional for general relativity is to
express the gravitational action in terms of \((h_{ab},N,N_a)\) and their time and
space derivatives. To simplify the discussion here, we will defer the analysis of
boundary terms until the end of this section. Thus, we start with the Hilbert action
\eqref{eq:E.1.13} rather than \eqref{eq:E.1.42} and for the present will discard the
boundary terms which arise in subsequent calculations. We express the scalar
curvature, \(R\), as
\begin{equation}
  R=2(G_{ab}n^an^b-R_{ab}n^an^b).
  \tag{E.2.26}\label{eq:E.2.26}
\end{equation}
From Eq.~\eqref{eq:10.2.30} we have
\begin{equation}
  G_{ab}n^an^b
    =\frac{1}{2}\left[{}^{(3)}R-K_{ab}K^{ab}+K^2\right],
  \tag{E.2.27}\label{eq:E.2.27}
\end{equation}
where \(K_{ab}\) is the extrinsic curvature of \(\Sigma_t\) and \(K=K^a{}_a\). On
the other hand, from the definition of the Riemann tensor, we have
\begin{align}
  R_{ab}n^an^b
    &=R_{acb}{}^c n^an^b\notag\\
    &=-n^a(\nabla_a\nabla_c-\nabla_c\nabla_a)n^c\notag\\
    &=(\nabla_an^a)(\nabla_cn^c)-(\nabla_cn^a)(\nabla_an^c)
      -\nabla_a(n^a\nabla_cn^c)+\nabla_c(n^a\nabla_an^c)\notag\\
    &=K^2-K_{ac}K^{ac}-\nabla_a(n^a\nabla_cn^c)
      +\nabla_c(n^a\nabla_an^c).
  \tag{E.2.28}\label{eq:E.2.28}
\end{align}
The last two terms on the right-hand side of Eq.~\eqref{eq:E.2.28} are divergences
and thus will be discarded. Hence, from Eqs.~\eqref{eq:E.1.12} and
\eqref{eq:E.2.25}--\eqref{eq:E.2.28}, we obtain
\begin{equation}
  \mathcal L_G=\sqrt{h}N\left[{}^{(3)}R+K_{ab}K^{ab}-K^2\right].
  \tag{E.2.29}\label{eq:E.2.29}
\end{equation}
The extrinsic curvature, \(K_{ab}\), is related to the "time derivative,"
\(\dot h_{ab}=h_a{}^ch_b{}^d\mathcal L_t h_{cd}\), of \(h_{ab}\) by
\begin{align}
  K_{ab}
    &=\frac{1}{2}\mathcal L_nh_{ab}
      =\frac{1}{2}\bigl[n^c\nabla_ch_{ab}
        +h_{ac}\nabla_bn^c+h_{cb}\nabla_an^c\bigr]\notag\\
    &=\frac{1}{2}N^{-1}\bigl[Nn^c\nabla_ch_{ab}
        +h_{ac}\nabla_b(Nn^c)+h_{cb}\nabla_a(Nn^c)\bigr]\notag\\
    &=\frac{1}{2}N^{-1}h_a{}^ch_b{}^d
      [\mathcal L_t h_{cd}-\mathcal L_Nh_{cd}]\notag\\
    &=\frac{1}{2}N^{-1}\bigl[\dot h_{ab}-D_aN_b-D_bN_a\bigr],
  \tag{E.2.30}\label{eq:E.2.30}
\end{align}
where \(D_a\) is the derivative operator on \(\Sigma_t\) associated with \(h_{ab}\)
(see chapter 10) and Eq.~\eqref{eq:E.2.23} was used to go from the second line to
the third line. Thus, substitution of Eq.~\eqref{eq:E.2.30} into
Eq.~\eqref{eq:E.2.29} expresses the gravitational action in the desired form given
by Arnowitt, Deser, and Misner (1962). The momentum canonically conjugate to
\(h_{ab}\) is
\begin{equation}
  \pi^{ab}=\frac{\partial\mathcal L_G}{\partial\dot h_{ab}}
    =\sqrt{h}(K^{ab}-Kh^{ab}).
  \tag{E.2.31}\label{eq:E.2.31}
\end{equation}

However, \(\mathcal L_G\) does not contain any time derivatives of \(N\) or \(N_a\),
so their conjugate momenta vanish identically. In analogy with the electromagnetic
case, we interpret this fact as telling us that \(N\) and \(N_a\) should not be
viewed as dynamical variables. Hence, we redefine our configuration space to consist
of Riemannian metrics, \(h_{ab}\), on \(\Sigma_t\). We define our Hamiltonian
density by
\begin{align}
  \mathcal H_G
    &=\pi^{ab}\dot h_{ab}-\mathcal L_G\notag\\
    &=-h^{1/2}N{}^{(3)}R
      +Nh^{-1/2}\left[\pi^{ab}\pi_{ab}-\frac{1}{2}\pi^2\right]
      +2\pi^{ab}D_aN_b\notag\\
    &=h^{1/2}\Bigl\{
      N\left[-{}^{(3)}R+h^{-1}\pi^{ab}\pi_{ab}
        -\frac{1}{2}h^{-1}\pi^2\right]
      -2N_b\left[D_a(h^{-1/2}\pi^{ab})\right]\notag\\
    &\hspace{5em}+2D_a(h^{-1/2}N_b\pi^{ab})\Bigr\},
  \tag{E.2.32}\label{eq:E.2.32}
\end{align}
where \(\pi=\pi^a{}_a\). Again, the last term in Eq.~\eqref{eq:E.2.32} contributes
only a boundary term to \(H_G=\int\mathcal H_G{}^{(3)}e\) and will be dropped.
Variation of \(H_G\) with respect to \(N\) and \(N_a\) yields the equations
\begin{equation}
  -{}^{(3)}R+h^{-1}\pi^{ab}\pi_{ab}
    -\frac{1}{2}h^{-1}\pi^2=0,
  \tag{E.2.33}\label{eq:E.2.33}
\end{equation}
\begin{equation}
  D_a(h^{-1/2}\pi^{ab})=0,
  \tag{E.2.34}\label{eq:E.2.34}
\end{equation}
which, with the substitution \eqref{eq:E.2.31}, can be recognized as the initial
value constraint equations \eqref{eq:10.2.28} and \eqref{eq:10.2.30} found in
chapter 10. The dynamical equations \eqref{eq:E.2.2} and \eqref{eq:E.2.3}
obtained from \(H_G\) are (Arnowitt, Deser, and Misner 1962)
\begin{equation}
  \dot h_{ab}=\frac{\delta H_G}{\delta\pi^{ab}}
    =2h^{-1/2}N\left(\pi_{ab}-\frac{1}{2}h_{ab}\pi\right)
      +2D_{(a}N_{b)},
  \tag{E.2.35}\label{eq:E.2.35}
\end{equation}
\begin{align}
  \dot\pi^{ab}
    &=-\frac{\delta H_G}{\delta h_{ab}}
      =-Nh^{1/2}\left({}^{(3)}R^{ab}
        -\frac{1}{2}{}^{(3)}Rh^{ab}\right)\notag\\
    &\quad+\frac{1}{2}Nh^{-1/2}h^{ab}
      \left(\pi_{cd}\pi^{cd}-\frac{1}{2}\pi^2\right)\notag\\
    &\quad-2Nh^{-1/2}
      \left(\pi^{ac}\pi_c{}^b-\frac{1}{2}\pi\pi^{ab}\right)\notag\\
    &\quad+h^{1/2}\left(D^aD^bN-h^{ab}D^cD_cN\right)\notag\\
    &\quad+h^{1/2}D_c(h^{-1/2}N^c\pi^{ab})
      -2\pi^{c(a}D_cN^{b)},
  \tag{E.2.36}\label{eq:E.2.36}
\end{align}
where, again, boundary terms have been ignored and Eq.~\eqref{eq:E.2.34} was used.
Equations \eqref{eq:E.2.33}--\eqref{eq:E.2.36} are equivalent to the vacuum
Einstein equation, \(R_{ab}=0\). Thus, we have succeeded in giving a constrained
Hamiltonian formulation of Einstein's equation.

The presence of constraints in our Hamiltonian formulations of Maxwell's equations
and Einstein's equations indicates that we have not isolated the "true dynamical
degrees of freedom" in our choice of configuration space. Even though we already
have eliminated \(V\) and \(N\) and \(N_a\) as dynamical variables, the constraints
tell us that our phase space still is "too large." This, in turn, is directly related
to the gauge freedom present in our configuration variables \(\mathbf{A}\) and
\(h_{ab}\), respectively. In the Maxwell case, \(\mathbf{A}\) and
\(\mathbf{A}-\boldsymbol{\nabla}\chi\) represent the same physical configuration. This suggests
that we should take our configuration space to be not simply the space of vector
potentials, \(\mathbf{A}\), but the space of equivalence classes, \(\overline{\mathbf{A}}\),
of vector potentials, where two vector potentials are equivalent if they differ only
by a gauge transformation. The cotangent space at \(\overline{\mathbf{A}}\) then would
be the space of linear functions of variations of \(\mathbf{A}\) which depend only on
the equivalence class. Thus, the momenta would be represented by vector fields
\(\boldsymbol{\pi}\) having the property that
\begin{equation}
  \int\boldsymbol{\pi}\cdot[\delta\mathbf{A}-\boldsymbol{\nabla}(\delta\chi)]
    =\int\boldsymbol{\pi}\cdot\delta\mathbf{A}.
  \tag{E.2.37}\label{eq:E.2.37}
\end{equation}
However, this property holds if and only if
\begin{equation}
  \boldsymbol{\nabla}\cdot\boldsymbol{\pi}=0.
  \tag{E.2.38}\label{eq:E.2.38}
\end{equation}
Thus, with our new choice of configuration space, the momentum space consists
precisely of the divergence-free vector fields on \(\Sigma_t\). But this means that
the constraint \eqref{eq:E.2.18} is automatically satisfied! We may drop the term
\(V(\boldsymbol{\nabla}\cdot\boldsymbol{\pi})\) from Eq.~\eqref{eq:E.2.16} and take the Hamiltonian
density to be simply
\begin{equation}
  \mathcal H_{\mathrm{EM}}
    =\frac{1}{2}(\boldsymbol{\pi}\cdot\boldsymbol{\pi}+\mathbf{B}\cdot\mathbf{B}),
  \tag{E.2.39}\label{eq:E.2.39}
\end{equation}
where, again, \(\mathbf{B}\equiv\boldsymbol{\nabla}\times\mathbf{A}\). (Note that \(\mathbf{B}\)
depends only on the equivalence class of \(\mathbf{A}\) and hence is a well defined
function of \(\overline{\mathbf{A}}\).) Hamilton's equations of motion
\eqref{eq:E.2.2} and \eqref{eq:E.2.3} yield
\begin{equation}
  \dot{\overline{\mathbf{A}}}
    =\frac{\delta\overline H_{\mathrm{EM}}}{\delta\boldsymbol{\pi}}=\boldsymbol{\pi},
  \tag{E.2.40}\label{eq:E.2.40}
\end{equation}
\begin{equation}
  \dot{\boldsymbol{\pi}}
    =-\frac{\delta\overline H_{\mathrm{EM}}}{\delta\overline{\mathbf{A}}}
    =-\boldsymbol{\nabla}\times\mathbf{B}.
  \tag{E.2.41}\label{eq:E.2.41}
\end{equation}
The equivalence classes appearing on both sides of Eq.~\eqref{eq:E.2.40} can be
eliminated by taking the curl of this equation. It then easily may be verified that
Eqs.~\eqref{eq:E.2.40} and \eqref{eq:E.2.41} with
\(\mathbf{E}\equiv-\boldsymbol{\pi}\) are equivalent to Maxwell's equations, where we remind
the reader that \(\boldsymbol{\nabla}\cdot\mathbf{B}=0\) follows automatically from the
definition of \(\mathbf{B}\), whereas \(\boldsymbol{\nabla}\cdot\mathbf{E}=0\) follows automatically
from our construction of the space of momenta. Thus, by eliminating the gauge
degrees of freedom in our configuration space by working with
\(\overline{\mathbf{A}}\) rather than \(\mathbf{A}\), we have succeeded in giving a
constraint-free Hamiltonian formulation of Maxwell's equations.\footnote{The
relation between the constraint \(\boldsymbol{\nabla}\cdot\mathbf{E}=0\) and the gauge
transformations \(\mathbf{A}\mapsto\mathbf{A}-\boldsymbol{\nabla}\chi\) can be obtained more
systematically as follows. Given a function \(f\) on phase space, we may associate
with it a vector field \(V\) on phase space by the requirement that for any function
\(g\) on phase space, we have \(V(g)=\{f,g\}\), where the \emph{Poisson bracket}
\(\{f,g\}\) is defined by
\[
  \{f,g\}=\int_{\Sigma_t}\left(
    \frac{\delta f}{\delta q}\frac{\delta g}{\delta\pi}
    -\frac{\delta g}{\delta q}\frac{\delta f}{\delta\pi}\right).
\]
(We have not defined infinite-dimensional manifolds here or vector fields on them,
so these remarks are intended only as heuristic.) One may verify directly that the
vector field \(V\) associated in this manner with the "constraint function"
\(f=\int_{\Sigma_t}\chi\boldsymbol{\nabla}\cdot\mathbf{E}\) (where \(\chi\) is an arbitrary
function on \(\Sigma_t\)) is just the infinitesimal generator of the one-parameter
family of transformations on phase space associated with the gauge transformations
\(\mathbf{A}\mapsto\mathbf{A}-\boldsymbol{\nabla}\chi\). Thus, in this sense, in electromagnetism,
the constraint "generates" the gauge transformations. By restricting to the
"constraint submanifold" \(\boldsymbol{\nabla}\cdot\mathbf{E}=0\) and to the space of orbits of
\(V\) on this submanifold, we obtain a consistent, constraint-free Hamiltonian
formulation on a "reduced phase space." Similarly, in the gravitational case, the
vector field associated with the constraint function
\(2h^{1/2}\int_{\Sigma_t}\xi^aD_a(h^{-1/2}\pi^b{}_a)\), where \(\xi^a\) is an
arbitrary vector field on \(\Sigma_t\), generates the one-parameter family of
diffeomorphisms on \(\Sigma_t\) associated with \(\xi^a\). Thus, one is led to
choose as the new configuration space the metrics on \(\Sigma_t\) modulo
diffeomorphisms (or, more precisely, modulo diffeomorphisms which can be
continuously deformed to the identity; see Friedman and Sorkin 1980).}

Similarly, in the case of Einstein's equation there is gauge arbitrariness in our
choice of configuration field \(h_{ab}\). If \(\psi\) is any diffeomorphism of
\(\Sigma_t\), then \(h_{ab}\) and \(\psi^*h_{ab}\) represent the same physical
configuration. This suggests that we should take the configuration space of general
relativity to be the set equivalence classes, \(\bar h_{ab}\), of Riemannian metrics
on \(\Sigma_t\), where two metrics are considered equivalent if they can be carried
into each other by a diffeomorphism. This configuration space is known as
\emph{superspace} (Wheeler 1968). Using superspace as the configuration space, we
find that for any vector field \(w^a\) on \(\Sigma_t\), the conjugate momenta
\(\pi^{ab}\) now must satisfy
\begin{equation}
  \int\pi^{ab}\bigl(\delta h_{ab}+D_{(a}w_{b)}\bigr)
    =\int\pi^{ab}\delta h_{ab},
  \tag{E.2.42}\label{eq:E.2.42}
\end{equation}
which implies that \(\pi^{ab}\) automatically satisfies
\begin{equation}
  D_a(h^{-1/2}\pi^{ab})=0.
  \tag{E.2.43}\label{eq:E.2.43}
\end{equation}
Thus, the constraint \eqref{eq:E.2.34} is eliminated by the choice of superspace as
the configuration space.

However, the constraint \eqref{eq:E.2.33} remains. This constraint may be viewed
as resulting from the gauge arbitrariness involved in the choice of how to "slice"
spacetime into space and time. It is very closely analogous to the constraint which
arises when one "parameterizes" an originally unconstrained theory in a fixed,
background spacetime, i.e., when one introduces into the Lagrangian a time
function---which defines the choice of hypersurfaces, \(\Sigma_t\), with respect to
a reference surface \(\Sigma\)---and treats this time function as a dynamical
variable (Kuchař 1973, 1981). In the case of such parameterized theories, the
constraint analogous to \eqref{eq:E.2.33} is linear in the momentum conjugate to
the time function. One then can "deparameterize" the theory by solving the
constraint for this momentum. However, in the case of Einstein's equation, the
constraint \eqref{eq:E.2.33} is quadratic in the momentum, and a similar
deparameterization does not appear to be possible. Thus, it does not appear possible
to find a choice of configuration space for general relativity such that only the
"true dynamical degrees of freedom" are present in its phase space. The presence of
the constraint \eqref{eq:E.2.33} appears to be an unavoidable feature of the
Hamiltonian formulation of general relativity. This provides a serious obstacle to
the formulation of a quantum theory of gravitation by the canonical quantization
approach (see chapter 14).

Finally, we return to the issue of boundary terms in the Hamiltonian formulation.
Consider, first, the case of a closed universe, i.e., \(M=\mathbb R\times\Sigma\),
where \(\Sigma\) is compact. Consider the region, \(U\), of \(M\) bounded by two
constant time hypersurfaces \(\Sigma_1\) and \(\Sigma_2\). Then the modified
gravitational action \(S'_G\), Eq.~\eqref{eq:E.1.42}, will receive boundary
contributions from \(\Sigma_1\) and \(\Sigma_2\). However, these boundary terms
will be canceled by the contributions of the third term on the right-hand side of
Eq.~\eqref{eq:E.2.28}. Note that the fourth term in Eq.~\eqref{eq:E.2.28} yields
no boundary contributions since \(n^a\nabla_an^c\) is orthogonal to the normal,
\(n^c\), to \(\Sigma_1\) and \(\Sigma_2\). In addition, the last term in
Eq.~\eqref{eq:E.2.32} makes no contribution since there is no spatial boundary.
Thus, in the case of a closed universe, our final answer for \(H_G\) is unchanged
when all boundary terms are reinserted. Note that because of
Eqs.~\eqref{eq:E.2.33} and \eqref{eq:E.2.34}, the numerical value of \(H_G\)
vanishes for every solution. This suggests that we should define the total energy of
a closed universe to be zero. In other words, it suggests that there does not exist a
nontrivial notion of total energy in a closed universe. However, this argument is not
conclusive, since one always can "parameterize" a theory in the manner mentioned
above so as to make its Hamiltonian vanish. If general relativity could be
"deparameterized," a notion of total energy in a closed universe could well emerge.

Consider, now, the case of asymptotically flat spacetimes. Again, consider a region
of \(M\) bounded by two hypersurfaces \(\Sigma_1\) and \(\Sigma_2\). As before, we
wish to consider metric variations for which \(h_{ab}\) is held fixed on
\(\Sigma_1\) and \(\Sigma_2\), but now the most natural spatial boundary condition
is that the variations preserve asymptotic flatness rather than that the induced
metric be held fixed on a distant spatial boundary. This new boundary condition
requires the addition of further boundary terms into the gravitational action
\eqref{eq:E.1.42}. Furthermore, the boundary terms\footnote{To avoid confusion, we
remind the reader that in Eq.~\eqref{eq:E.1.42} \(K\) is the trace of the
extrinsic curvature of the boundary, whereas in Eq.~\eqref{eq:E.2.28} \(K_{ab}\)
is the extrinsic curvature of \(\Sigma_t\).} from Eqs.~\eqref{eq:E.2.28} and
\eqref{eq:E.2.32} now will contribute to \(H_G\). Instead of keeping careful
account of all these terms, we shall proceed by calculating the boundary terms
arising from variations of \(H_G\) and then modifying \(H_G\) to get rid of these
terms. Introduce on \(\Sigma_t\) asymptotic Cartesian coordinates as described in
problem 2 of chapter 11. We consider the case where \(t^a\) asymptotically becomes
a time translation at spatial infinity, i.e., we take \(N\to1\) and \(N_a\to0\) as
\(r\to\infty\). Let \(S\) denote a coordinate sphere of radius \(r\). Then, only
the term \(-h^{1/2}N{}^{(3)}R\) in Eq.~\eqref{eq:E.2.32} produces a nonvanishing
boundary term on \(S\) in the limit \(r\to\infty\) when \(H_G\) is varied. But the
calculation of this term is the three-dimensional analog of the calculation of the
boundary term in the gravitational action given at the end of the previous section,
except that we no longer require the metric to be held fixed at \(S\). Thus, using
Eq.~\eqref{eq:7.5.14} (or, even better, using the three-dimensional analog of the
first line of Eq.~\eqref{eq:E.1.38}) we find that for a one-parameter variation of
\(h_{ab}\) and \(\pi^{ab}\) which preserves asymptotic flatness, we have (Regge and
Teitelboim 1974)
\begin{equation}
  \frac{dH_G}{d\lambda}
    =\int_{\Sigma_t}\left[A_{ab}\delta\pi^{ab}-B^{ab}\delta h_{ab}\right]
      -\delta C.
  \tag{E.2.44}\label{eq:E.2.44}
\end{equation}
Here \(A_{ab}\) and \(B^{ab}\) are given by the right-hand sides of
Eqs.~\eqref{eq:E.2.35} and \eqref{eq:E.2.36}, respectively, and the boundary term
\(\delta C\) is given by
\begin{equation}
  \delta C=\lim_{r\to\infty}\int_S
    r^ah^{bc}\left[D_c(\delta h_{ab})-D_a(\delta h_{bc})\right],
  \tag{E.2.45}\label{eq:E.2.45}
\end{equation}
where \(r^a\) is the unit normal to \(S\) and the natural volume element on \(S\)
is understood. We can rewrite \(\delta C\) in coordinate component form as
\begin{align}
  \delta C
    &=\lim_{r\to\infty}\sum_{\nu=1}^{3}\int_S
      \left(\frac{\partial\delta h_{\mu\nu}}{\partial x^\nu}
        -\frac{\partial\delta h_{\nu\nu}}{\partial x^\mu}\right)r^\mu\notag\\
    &=\delta\left\{\lim_{r\to\infty}\sum_{\nu=1}^{3}\int_S
      \left(\frac{\partial h_{\mu\nu}}{\partial x^\nu}
        -\frac{\partial h_{\nu\nu}}{\partial x^\mu}\right)r^\mu\right\},
  \tag{E.2.46}\label{eq:E.2.46}
\end{align}
where we have discarded terms which do not contribute in the limit as
\(r\to\infty\). Thus, in order to get a Hamiltonian whose variation produces no
boundary terms from spatial infinity, we define a new gravitational Hamiltonian
\(H'_G\) by
\begin{equation}
  H'_G=H_G+\alpha,
  \tag{E.2.47}\label{eq:E.2.47}
\end{equation}
where \(\alpha\) is the term inside the braces in Eq.~\eqref{eq:E.2.46}. The
right-hand sides of Eqs.~\eqref{eq:E.2.35} and \eqref{eq:E.2.36} then truly are the
functional derivatives of \(H'_G\) with respect to \(h_{ab}\) and \(\pi^{ab}\) for
variations which preserve asymptotic flatness.

The numerical value of \(H'_G\) for a solution of Einstein's equation is just
\(\alpha\). This suggests that \(\alpha\) should be interpreted as proportional to
the total energy of an asymptotically flat spacetime. This provides the motive for
the definition given in chapter 11, Eq.~\eqref{eq:11.2.14}. (The constant of
proportionality between \(\alpha\) and energy can be determined by evaluating
\(\alpha\) for the Schwarzschild solution.) Similarly, the definition of total
momentum, Eq.~\eqref{eq:11.2.15}, can be motivated by examining the boundary
terms in \(H_G\) which occur when we take \(N\to0\) and require \(N_a\) to go to a
translation as \(r\to\infty\). Indeed, a notion of angular momentum (see problem 6
of chapter 11) arises from consideration of more general asymptotic behavior of the
lapse and shift (Regge and Teitelboim 1974).
